\begin{figure}[t]
\centering
\resizebox{\linewidth}{!}{\begin{tikzpicture}[node distance=5mm and 5mm]
  \node[box] (x) {state $s_t$\\$C$ fields};
  \node[box, right=of x] (p) {perceive\\$3{\times}3$ circular\\(dilated in -MS)};
  \node[box, right=of p] (m) {shared\\per-cell MLP};
  \node[box, right=of m] (f) {flux head\\$2C$ channels\\zero-init};
  \node[box, right=of f] (d) {per-field\\divergence $D_h$};
  \node[op, right=5mm of d] (add) {$+$};
  \node[box, right=of add] (b) {admissible-state\\projection\\(-B only)};
  \node[box, right=of b] (o) {$s_{t+1}$};
  \draw[flow] (x) -- (p);
  \draw[flow] (p) -- (m);
  \draw[flow] (m) -- (f);
  \draw[flow] (f) -- (d);
  \draw[flow] (d) -- (add);
  \draw[flow] (add) -- (b);
  \draw[flow] (b) -- (o);
  \draw[flow] (x.south) to[out=-35,in=-145] (add.south);
  \node[font=\scriptsize, below=7mm of d, align=center] (n1)
    {$\sum_{\mathrm{cells}} (D_hF)_c = 0$ for every field $c$\\and every $\theta$ (Prop.~\ref{prop:conservation})};
  \node[font=\scriptsize, below=7mm of b, align=center] (n2)
    {clip to $[\ell,h]$, then restore mass\\in proportion to headroom (Prop.~\ref{prop:headroom})};
  \draw[dashed,gray] (d.south) -- (n1.north);
  \draw[dashed,gray] (b.south) -- (n2.north);
\end{tikzpicture}}
\caption{One PI-NCA step. A single local rule, shared by every cell, maps the
neighbourhood of each cell to a flux for each field rather than to a state increment.
The periodic discrete divergence of that flux cannot change any field's total, whatever
the weights are, so conservation is a property of the architecture and not of training.
MC-PI-NCA is this construction for $C>1$; the suffix -MS widens perception with
dilations $(1,2,4)$ and -B appends the bound-preserving projection of
Equation~\ref{eq:headroom}.}
\label{fig:pinca}
\end{figure}
